\documentclass[10pt]{article}

\usepackage[margin=0.75in]{geometry}
\usepackage{amsmath,amsthm,amssymb,bm}
\usepackage{xcolor}
\usepackage{cancel}
\usepackage{graphicx}
\usepackage{changepage}
\usepackage{circuitikz}
\usepackage{pgfplots}
\usepackage{minted}
\usepackage{physics}
\usepackage{tikz}
\usepackage{tikz-3dplot}
\usepackage{hyperref}
\usepackage{cleveref}
\usepackage{siunitx}
\usepackage{fontspec}
\usepackage{relsize}
\usepackage{subfig}
\usepackage{todonotes}
\usepackage{contour}
\usepackage{pgffor}
\usepackage{multicol, multirow, booktabs}
\usepackage[breakable]{tcolorbox}
\usepackage[inline]{enumitem}

\theoremstyle{definition}
\newtheorem{problem}{Problem}
\newtheorem{soln}{Solution}

\pgfplotsset{compat=newest}
\usetikzlibrary{lindenmayersystems}
\usetikzlibrary{arrows}
\usetikzlibrary{calc}
\usetikzlibrary{positioning, fit}
\usetikzlibrary{3d, perspective}
\usetikzlibrary{math}
\usetikzlibrary{fadings}
\usetikzlibrary{angles,quotes} 
\usetikzlibrary{decorations.markings}

\definecolor{incolor}{HTML}{303F9F}
\definecolor{outcolor}{HTML}{D84315}
\definecolor{cellborder}{HTML}{CFCFCF}
\definecolor{cellbackground}{HTML}{F7F7F7}
\newcommand{\ui}{\hat{i}}
\newcommand{\uj}{\hat{j}}
\newcommand{\uk}{\hat{k}}
\newcommand{\ux}{\hat{\mathbf{x}}}
\newcommand{\uy}{\hat{\mathbf{y}}}
\newcommand{\uz}{\hat{\mathbf{z}}}
\newcommand{\uv}[1]{\hat{\bv{{#1}}}}
\newcommand{\pr}[1]{{#1^\prime}}
\newcommand{\pd}[2][]{{\frac{\partial^{#1}}{\partial {{#2}^{#1}}}}}
\newcommand{\dif}[2][]{{\frac{d^{#1}}{d{{#2}^{#1}}}}}
\newcommand{\grd}{\nabla}

\pgfdeclarelayer{background}  
\pgfsetlayers{background,main}
\AtBeginDocument{\RenewCommandCopy\qty\SI}

\makeatletter
\newcommand{\boxspacing}{\kern\kvtcb@left@rule\kern\kvtcb@boxsep}
\makeatother
\newcommand{\prompt}[4]{
    \ttfamily\llap{{\color{#2}[#3]:\hspace{3pt}#4}}\vspace{-\baselineskip}
}

\newcommand{\thevenin}[2]{
  \begin{center}
    \begin{circuitikz} \draw
      (0,0) -- (2,0) to[battery1, l_=$V_{Th}\eq#1$] (2,2) 
      to[resistor, l_=$R_{Th}\eq#2$] (0,2)
      ;
      \draw [o-] (-.07,2.079);
      \draw [o-] (-.07,0.079);
    \end{circuitikz}
  \end{center}
}

\newcommand{\norton}[2]{
  \begin{center}
    \begin{circuitikz} \draw
      (0,0) -- (3,0) to[american current source, l_=$I_{N}\eq#1$] (3,2) -- (0,2) (2,0)
      to[resistor, l=$R_{N}\eq#2$] (2,2)
      ;
      \draw [o-] (-.07,2.079);
      \draw [o-] (-.07,0.079);
    \end{circuitikz}
  \end{center}
}

\DeclareMathOperator{\Div}{div}
\DeclareMathOperator{\Curl}{curl}
\DeclareMathOperator{\sgn}{sgn}

\newcommand{\highlight}[1]{\colorbox{yellow}{$\displaystyle #1$}}
\newcommand{\ti}[1]{\widetilde{#1}}
\renewcommand{\div}[1]{\bv{\nabla}\cdot #1}
\renewcommand{\curl}[1]{\bv{\nabla}\times #1}

\newfontface{\Kaufmann}{Kaufmann}
\DeclareTextFontCommand{\kf}{\Kaufmann}
\newcommand{\scriptr}{\fontsize{12pt}{12pt}\kf{r}\fontsize{10pt}{10pt}}

\newfontface{\KaufmannB}{Kaufmann Bold}
\DeclareTextFontCommand{\kfb}{\KaufmannB}
\newcommand{\bscriptr}{\fontsize{12pt}{12pt}\kfb{r}\fontsize{10pt}{10pt}}
\newcommand{\justif}[2]{&{#1}&\text{#2}}
\newcommand{\bv}[1]{\bm{\mathrm{#1}}}

\title{Physics 4140H: Assignment II}
\author{Jeremy Favro (0805980) \\ Trent University, Peterborough, ON, Canada}
\date{\today}

\begin{document}
\maketitle

% PROBLEM 1
\begin{problem} By using the Euler equations with an auxiliary condition imposed, find the shortest path between
any two points on the surface of a sphere.
\end{problem}
\begin{soln} We impose the constraint that $g(x)=x^2+y^2+z^2-R^2=0$, i.e. that we are on the surface of a sphere of radius
  $R$. Euler's equations are
  \begin{align*}
    \frac{\partial f}{\partial y}-\frac{d}{dx}\left(\frac{\partial f}{\partial y'}\right)+\lambda(x)\frac{\partial g}{\partial y} & =0 \\
    \frac{\partial f}{\partial z}-\frac{d}{dx}\left(\frac{\partial f}{\partial z'}\right)+\lambda(x)\frac{\partial g}{\partial z} & =0
  \end{align*}
  where $f$ takes the usual form for distance-measure, $f=\sqrt{1+(y')^2+(z')^2}$. Hence, the first two derivatives are zero
  and
  $$
    f_{y'}=\frac{y'}{\sqrt{1+(y')^2+(z')^2}}
  $$
  and $$f_{z'}=\frac{z'}{\sqrt{1+(y')^2+(z')^2}}$$
  and $g_y=2y$ and $g_z=2z$. So, Euler's equations become
  \begin{align*}
    \frac{d}{dx}\left(\frac{y'}{\sqrt{1+(y')^2+(z')^2}}\right) & =2y\lambda(x) \\
    \frac{d}{dx}\left(\frac{z'}{\sqrt{1+(y')^2+(z')^2}}\right) & =2z\lambda(x)
  \end{align*}
  which we can rewrite as a single equation to arrive at
  $$\frac{1}{y}\frac{d}{dx}\left(\frac{y'}{\sqrt{1+(y')^2+(z')^2}}\right)=\frac{1}{z}\frac{d}{dx}\left(\frac{z'}{\sqrt{1+(y')^2+(z')^2}}\right)$$
  this is tantalizingly close to something nice but I can't think of a way to incorporate the $1/y$ and $1/z$ so I could get rid of the derivatives, so
  we'll take the derivatives in full detail. So,
  $$\frac{d}{dx}\left(\frac{y'}{\sqrt{1+(y')^2+(z')^2}}\right)=\frac{y''\sqrt{1+(y')^2+(z')^2}-\frac{y'(y'y''+z'z'')}{\sqrt{1+(y')^2+(z')^2}}}{1+(y')^2+(z')^2}$$
  and similarly
  $$\frac{d}{dx}\left(\frac{z'}{\sqrt{1+(y')^2+(z')^2}}\right)=\frac{z''\sqrt{1+(y')^2+(z')^2}-\frac{z'(y'y''+z'z'')}{\sqrt{1+(y')^2+(z')^2}}}{1+(y')^2+(z')^2}$$
  and so the equality becomes, multiplying by $yz$ on either side and cancelling the common denominator,
  $$zy''\sqrt{1+(y')^2+(z')^2}-\frac{zy'(y'y''+z'z'')}{\sqrt{1+(y')^2+(z')^2}}=yz''\sqrt{1+(y')^2+(z')^2}-\frac{yz'(y'y''+z'z'')}{\sqrt{1+(y')^2+(z')^2}}$$
  multiplying by the radical
  $$zy''+\cancel{z(y')^2y''}+z(z')^2y''-\cancel{z(y')^2y''}-zz'z''y'=z''y+y(y')^2z''+\cancel{(z')^2z''y}-z'yy'y''-\cancel{(z')^2z''y}$$
  so
  $$zy''+z(z')^2y''-zz'z''y'=z''y+y(y')^2z''-z'yy'y''$$
  now we want to obtain, using the hint, an expression involving $x$. The only way we can get that is by using the
  $g(x)$ condition. Note that
  $$\frac{dg}{dx}=x+yy'+zz'=0\text{ or } -x=yy'+zz'$$
  so let's factor the equality to try to obtain $yy'+zz'$. Rearranging,
  $$zy''+z(z')^2y''+z'yy'y''=z''y+y(y')^2z''+zz'z''y'$$
  and factoring out the second derivatives
  $$\left(z+z(z')^2+z'yy'\right)y''=\left(y+y(y')^2+zz'y'\right)z''$$
  and pulling out the extra $z'$ and $y'$ terms from the second two terms in each set of parentheses,
  $$\left(z+\left[zz'+yy'\right]z'\right)y''=\left(y+\left[yy'+zz'\right]y'\right)z''$$
  and so using the derivative of the constraint,
  $$\left(z-xz'\right)y''=\left(y-xy'\right)z''.$$
  Yay! Now to show that the equation of a plane passing through the origin, $Ax+By=z$, is a solution here.
  Note that $z'=A+By'$ and $z''=By''$. Then $y=\frac{1}{B}\left(z-Ax\right)$, $y'=\frac{1}{B}\left(z'-A\right)$ and $y''=\frac{1}{B}z''$ so
  plugging these in,
  $$\left(Ax+By-x(A+By')\right)\left(\frac{1}{B}z''\right)=\left(\frac{1}{B}\left(z-Ax\right)-x\left(\frac{1}{B}\left(z'-A\right)\right)\right)By''.$$
  Multiplying things out
  $$\left(\cancel{Ax}+By-\cancel{Ax}-Bxy'\right)\left(\frac{1}{B}z''\right)=\left(\frac{1}{B}z-\cancel{\frac{A}{B}x}-x\frac{1}{B}z'+\cancel{\frac{A}{B}x}\right)By''$$
  and clearly by cancelling the remaining $B$ factors
  $$\left(y-xy'\right)z''=\left(z-xz'\right)y'',$$
  our original equation, reversed. This is a true statement, so the plane must satisfy the equation. The other way to see this is by noting that,
  since the the plane is linear in $x$, the second derivatives of $y$ and $z$ are zero and so, trivially, $0=0$.
\end{soln}

% PROBLEM 2
\begin{problem} Find the shortest path between $(x, y, z)$ points $(0,-1,0)$ and $(0,1,0)$ on the conical surface $z=1-\sqrt{x^2+y^2}$.
What is the length of the path? (Hint: use cylindrical coordinates).
\end{problem}
\begin{soln}
\end{soln}

% PROBLEM 3
\begin{problem} Find the shortest path between $(x, y, z)$ points $(0,-1,0)$ and $(0,1,0)$ on the conical surface $z=1-\sqrt{x^2+y^2}$.
What is the length of the path? (Hint: use cylindrical coordinates).
\end{problem}
\begin{soln}
\end{soln}
\end{document}